% Generated draft: requires mathematical and source review.
\documentclass[11pt]{article}
\usepackage[margin=25mm]{geometry}
\usepackage{fontspec}
\setmainfont{Latin Modern Roman}
\usepackage{amsmath,amssymb,mathtools,unicode-math}
\setmathfont{Latin Modern Math}
\usepackage{enumitem,needspace,xurl,hyperref}
\setlength{\parindent}{0pt}
\setlength{\parskip}{5pt}
\setlength{\emergencystretch}{3em}
\newcommand{\N}{\mathbb{N}}
\newcommand{\Z}{\mathbb{Z}}
\newcommand{\Q}{\mathbb{Q}}
\newcommand{\R}{\mathbb{R}}
\newcommand{\C}{\mathbb{C}}
\newcommand{\F}{\mathbb{F}}
\newcommand{\A}{\mathbb{A}}
\newcommand{\PP}{\mathbb{P}}
\newcommand{\E}{\mathbb{E}}
\newcommand{\eps}{\varepsilon}
\newcommand{\dd}{\,\mathrm{d}}
\begin{document}
\section*{UnsolvedMath ID: 30004479}
\textbf{Source problem code:} catalogue entry
\subsection{Definitions and mathematical statement}
Let \(\mathcal E\) denote the high-dimensional embedding space intended by the source, and let \(\pi_0(\mathcal E)\) denote its set of path components. Let \(\mathcal G\) denote the associated hairy graph complex. The source describes a natural finite-to-one map, written schematically as
\[
\Phi\colon \pi_0(\mathcal E)\longrightarrow \operatorname{MC}(\mathcal G),
\]
where \(\operatorname{MC}(\mathcal G)\) stands for the intended Maurer--Cartan target. The finite-to-one qualification means that each fiber \(\Phi^{-1}(\{\alpha\})\) is finite; it does not assert surjectivity. This notation is provisional: the supplied record does not specify the embedding space, the graph complex or its algebraic structure, or whether the target consists of individual Maurer--Cartan elements or equivalence classes.

\textbf{Source question.} What relationship, if any, does this map have to classical embedding invariants, including the Boechat--Haefliger invariant, the Whitney--Skopenkov invariant, and linking invariants, in settings where those invariants are defined?

\textbf{Unresolved formulation.} The source supplies neither the manifolds and dimension restrictions defining \(\mathcal E\), nor the hypotheses and conventions needed to define \(\mathcal G\), \(\Phi\), and the classical invariants on a common domain. It also does not specify what mathematical relationship is sought. Consequently, this is a schematic research question rather than a fully specified conjectural equality, factorization, or classification statement.
\subsection{Short English statement}
The question asks how a stated finite-to-one graph-complex invariant of embedding components relates to the Boechat--Haefliger, Whitney--Skopenkov, and linking invariants. The embedding spaces, Maurer--Cartan target, applicable hypotheses, and intended relationship remain unspecified.
\subsection{Sources}
\begin{itemize}
\item Supplied local UnsolvedMath registry record, permanent ID \(30004479\), source problem code \texttt{catalogue entry}, titled ``Graph Complex Invariants and Classical Embedding Invariants,'' in the Oberwolfach Reports Open Problems collection. No individual proposer is supplied; the proposed year is recorded as 2020.
\item Supplied report citation: \emph{Topologie} (2021), Oberwolfach Reports, pp.\ 1083--1137, DOI \texttt{10.4171/owr/2020/21}. \url{https://doi.org/10.4171/owr/2020/21}. The background describes these reports as documenting mathematical workshops and their open-problem discussions.
\item Supplied literature source: ``OWR Topology report.'' \url{https://ems.press/content/serial-article-files/46858}.
\item Supplied literature source: Crowley--Skopenkov, ``Embeddings of non-simply-connected 4-manifolds in 7-space.'' \url{https://www.mathnet.ru/eng/mmj787}.
\item Supplied literature assessment, checked \texttt{2026-08-22}: the record labels the question open and describes a conjectured relationship between a rational homotopy invariant of embedding spaces and classical embedding invariants. It states that later classifications of several embedding families do not establish the proposed general identification. This is supplied dated curation derived from the OWR record, not an independently verified literature assessment; the background explicitly calls for independent verification of time-sensitive status and cited claims.
\end{itemize}
\end{document}
